% Einheit 4 von 5 – Trapez, Drachenviereck, Raute (bank/flaechen/e4.jsonl, zusammenbau v0.1)
%% TODO Zweigzeile Teil 1: was hier gelernt wird (Satz in Schülersprache aus dem Einheitstitel) – Einheit 4
\einheitenkopf[e4]{Einheit 4 von 5 · Trapez, Drachenviereck, Raute}
\zweigzeile{ab Kl. 7, je nach Buch bis Kl. 8 · P10}
\setcounter{aufgabe}{21}

%% TODO Titel: Ich-kann-Satz fehlt in der Bank (Platzhalter: Kettenname) – Nr. 22
%% TODO Anweisung: ein Satz über der ersten Teilaufgabe fehlt in der Bank – Nr. 22
\begin{aufgabe}{Trapez}
%% TODO \rechenplatz{n}: Schrittzahl des Grundfalls fehlt in der Bank (nur bei fester Schreibform, 2.3 b)
\begin{teile}
\teil Welche Formel passt zur Figur? ($a$, $c$ parallele Seiten, $h$ Höhe, $e$, $f$ Diagonalen) \\ \kreuz{$A = a \cdot h$} \\ \kreuz{$A = \frac{a + c}{2} \cdot h$} \\ \kreuz{$A = \frac{e \cdot f}{2}$}

\trapez[hoehe=h]{5}{3}{2.5}
\teil Trapez, $a = 9$ cm, $c = 5$ cm, $h = 4$ cm – Fläche? A = \leerfeld[cm²]
\teil Trapez, $a = 7{,}5$ cm, $c = 4{,}5$ cm, $h = 3$ cm – Fläche? A = \leerfeld[cm²]
\teil Trapez, $A = 42$ cm², $a = 9$ cm, $c = 5$ cm – Höhe? h = \leerfeld[cm]
\teil Drachenviereck mit den Diagonalen $e = 7$ cm und $f = 6$ cm – Fläche? A = \leerfeld[cm²]
\teil Gleichschenkliges Trapez: $a = 11$ cm, $c = 5$ cm, Schenkel je $5$ cm, Höhe $4$ cm – Umfang? u = \leerfeld[cm]
\teil Ein Regalbrett hat die Form eines gleichschenkligen Trapezes mit den parallelen Seiten $90$ cm und $70$ cm; seine Fläche ist $2\,960$ cm². Passen Blumentöpfe mit $35$ cm Durchmesser darauf? Weise es über die Tiefe (Höhe) nach. \hfill (P10 2014 OS)
\end{teile}
\end{aufgabe}

%% TODO Titel: Ich-kann-Satz fehlt in der Bank (Platzhalter: Kettenname) – Nr. 23
%% TODO Anweisung: ein Satz über der ersten Teilaufgabe fehlt in der Bank – Nr. 23
\begin{aufgabe}{Umfang mit Schenkel aus Pythagoras oder Sinus (Vorrat)}
\begin{teile}
\teil Gleichschenkliges Trapez: $a = 14$ cm, $c = 8$ cm, $h = 4$ cm – Umfang? u = \leerfeld[cm]
\end{teile}
\end{aufgabe}

%% TODO Titel: Ich-kann-Satz fehlt in der Bank (Platzhalter: Kettenname) – Nr. 24
%% TODO Anweisung: ein Satz über der ersten Teilaufgabe fehlt in der Bank – Nr. 24
\begin{aufgabe}{Diagonale aus A}
\begin{teile}
\teil Drachenviereck, $A = 35$ cm², $e = 7$ cm – wie lang ist $f$? f = \leerfeld[cm]
\end{teile}
\end{aufgabe}

%% TODO Titel: Ich-kann-Satz fehlt in der Bank (Platzhalter: Kettenname) – Nr. 25
%% TODO Anweisung: ein Satz über der ersten Teilaufgabe fehlt in der Bank – Nr. 25
\begin{aufgabe}{Trapez – Fehler finden, Begründen, Anwendung}
\begin{teile}
\teil Trapez mit $a = 15$ cm, $c = 9$ cm, $h = 7$ cm. Leo rechnet: \rechnung{A &= 15 \cdot 7 \\ A &= 105 \text{ cm}^2} Wo ist der Fehler?
\teil Warum rechnet man beim Trapez „Mittelwert der parallelen Seiten mal Höhe“? Begründe mit zwei gleichen Trapezen.
\teil Ein Deich hat im Querschnitt die Form eines Trapezes: unten $32$ m, oben $6$ m breit, $7$ m hoch. Wie viel m³ Erde braucht man für $100$ m Deich?
\end{teile}
\end{aufgabe}
